% Worked Examples: Concept 2.5.4 - Block Diagrams and Transfer Functions
\documentclass[11pt,a4paper]{article}
\usepackage{worked-examples-template}

\title{\textbf{Worked Examples}\\[0.5em]
\large Concept 2.5.4: Block Diagrams and Transfer Functions\\[0.3em]
\normalsize \coursesubtitle{} -- \coursetitle}
\author{Assoc.\ Prof.\ William Robertson \and Dr Sean McGowan}
\date{\today}

\begin{document}

\maketitle
\tableofcontents
\newpage

\section{Concept 2.5.4: Block Diagrams}

\subsection{Example 1: Block Diagram Reduction}

\paragraph{Problem:} Consider a feedback control system with the following components in series and feedback configuration:
\begin{itemize}
\item Controller: $C(s) = K_p + \frac{K_i}{s}$ (PI controller)
\item Plant: $P(s) = \frac{1}{s+1}$
\item Feedback sensor: $H(s) = 1$ (unity feedback)
\end{itemize}

Find the closed-loop transfer function from reference $R(s)$ to output $Y(s)$.

\paragraph{Solution:}

For a standard feedback configuration with controller $C(s)$, plant $P(s)$, and feedback $H(s) = 1$, the closed-loop transfer function is:
\begin{align}
T(s) = \frac{C(s)P(s)}{1 + C(s)P(s)H(s)} = \frac{C(s)P(s)}{1 + C(s)P(s)}
\end{align}

First, find the product $C(s)P(s)$:
\begin{align}
C(s)P(s) &= \left(K_p + \frac{K_i}{s}\right) \cdot \frac{1}{s+1} \\
&= \frac{K_ps + K_i}{s} \cdot \frac{1}{s+1} \\
&= \frac{K_ps + K_i}{s(s+1)}
\end{align}

Now compute the closed-loop transfer function:
\begin{align}
T(s) &= \frac{\frac{K_ps + K_i}{s(s+1)}}{1 + \frac{K_ps + K_i}{s(s+1)}} \\
&= \frac{K_ps + K_i}{s(s+1) + K_ps + K_i} \\
&= \frac{K_ps + K_i}{s^2 + s + K_ps + K_i} \\
&= \frac{K_ps + K_i}{s^2 + (1+K_p)s + K_i}
\end{align}

\textbf{Result:} The closed-loop transfer function is:
\begin{align}
T(s) = \frac{K_ps + K_i}{s^2 + (1+K_p)s + K_i}
\end{align}

\textbf{Analysis:}
\begin{itemize}
\item The system is second-order with one zero at $s = -K_i/K_p$
\item The closed-loop poles depend on values of $K_p$ and $K_i$
\item For a specific example with $K_p = 2$ and $K_i = 3$:
\begin{align}
T(s) = \frac{2s + 3}{s^2 + 3s + 3}
\end{align}
The poles are at $s = \frac{-3 \pm \sqrt{9-12}}{2} = \frac{-3 \pm j\sqrt{3}}{2} \approx -1.5 \pm j0.866$
\end{itemize}
\end{document}
